\section{Related Work}
\label{sec:related}

\noindent\textbf{Repository-level developer question answering.}
A recent line of benchmarks evaluates LLMs on questions about real codebases.
SWE-QA~\cite{sweqa2026} constructs repository-level questions from a taxonomy of
seed templates instantiated per repository with program analysis and an LLM,
pairing them with RAG-assisted, expert-revised answers and an LLM-as-judge
evaluation with human validation. Its questions are clean and controllable but
deliberately synthetic: no developer ever asked them. CoReQA~\cite{coreqa2025}
similarly targets repository-level comprehension with code-centric context.
StackRepoQA~\cite{stackrepoqa2025} moves toward authenticity by pairing Stack
Overflow questions with repositories, and reports a finding that motivates much
of our design: a substantial share of apparent model accuracy on public
developer threads reflects memorization of those threads rather than reasoning
over the repository. All three treat developer questions as a uniform
population, supply only code-internal context, and grade against answers whose
correctness is itself difficult to verify externally. SecDevQA differs on each
axis: its questions are verbatim-grounded in real maintainer exchanges, are
specifically security questions whose answers depend on artifacts \emph{outside}
the repository (advisories, vulnerability databases), carry externally
verifiable facts that anchor grading, and are evaluated under a time cap
that makes the memorization risk StackRepoQA observed measurable per item
(\S\ref{sec:contamination}) rather than merely acknowledged.

\noindent\textbf{Security benchmarks for LLMs and agents.}
Security-specific evaluations target tasks adjacent to ours but distinct from
it. SecVulEval~\cite{secvuleval2025} measures vulnerability \emph{detection} on
labelled code; CodeSecEval~\cite{codeseceval2024} measures secure code
\emph{generation}; SecureAgentBench~\cite{secureagentbench2025} evaluates
whether coding agents produce secure patches under repository tasks; and
SecQA~\cite{secqa2023} tests textbook security knowledge through
multiple-choice questions. These benchmarks share deterministic or
near-deterministic grading, but none evaluates the question-answering task
developers actually pose to maintainers---free-text, situation-specific, and
resolved by consulting heterogeneous artifacts---and none treats artifact
access as a controlled variable. SecDevQA is complementary: where
SecureAgentBench asks whether an agent can \emph{write} a secure change, we ask
whether it can \emph{answer} the security questions that precede and motivate
such changes, and which evidence it needs to do so.

\noindent\textbf{Developer security information behavior.}
That developers' security outcomes depend on the information sources they
consult is well established. Acar et al.~\cite{acar2016information} showed
experimentally that developers permitted to consult Stack Overflow produced
less secure code than those using official documentation, and that only a
minority of consulted threads contained secure advice; Fischer et
al.~\cite{fischer2017stackoverflow} traced insecure code snippets from such
threads into hundreds of thousands of deployed Android applications. This
literature measures the \emph{consequences} of answer quality but offers no
instrument for evaluating a new class of answerers. As LLMs and agents become
the first responders to developer security questions, SecDevQA provides that
instrument, and its category-to-artifact map (RQ1) extends this literature with
a quantitative account of \emph{which evidence} trustworthy answers rest on.

\noindent\textbf{Positioning.}
Table~\ref{tab:related} summarizes the comparison. The claim we make is
precise: no existing benchmark combines (a)~real security Q\&A mined from
developer forums, (b)~external advisory artifacts as a context variable,
(c)~fact-grounded, condition-aware grading, and (d)~systematic measurement of
which artifact types are necessary for which question categories. We do not
claim to be the first repository-level QA benchmark, nor the first security
evaluation of LLMs.

\begin{table}[t]
\caption{Comparison with related work.
  \textbf{Sec.}~=~security-specific corpus;
  \textbf{Adv.}~=~external advisory artifacts as context;
  \textbf{Grd.}~=~fact-grounded grading;
  \textbf{Abl.}~=~artifact-type ablation.
  \checkmark~full; P~partial; $\times$~none.}
\label{tab:related}
\centering
\footnotesize
\setlength{\tabcolsep}{4pt}
\begin{tabular}{l c c c c}
\toprule
\textbf{Paper} & \textbf{Sec.} & \textbf{Adv.} & \textbf{Grd.} & \textbf{Abl.} \\
\midrule
SWE-QA~\cite{sweqa2026}            & $\times$ & $\times$ & P          & $\times$ \\
StackRepoQA~\cite{stackrepoqa2025} & $\times$ & $\times$ & $\times$   & $\times$ \\
CoReQA~\cite{coreqa2025}           & $\times$ & $\times$ & P          & $\times$ \\
SecVulEval / CodeSecEval           & \checkmark & $\times$ & \checkmark & $\times$ \\
SecureAgentBench                   & \checkmark & $\times$ & \checkmark & $\times$ \\
SecQA                              & \checkmark & $\times$ & \checkmark & $\times$ \\
\midrule
\textbf{SecDevQA (ours)}           & \checkmark & \checkmark & \checkmark & \checkmark \\
\bottomrule
\end{tabular}
\end{table}
